\documentclass[11pt]{article}
\usepackage[margin=1.05in]{geometry}
\usepackage{amsmath,amssymb,amsthm}
\usepackage{enumitem}
\usepackage{booktabs}

\newtheorem{lemma}{Lemma}[section]
\newtheorem{proposition}[lemma]{Proposition}
\newtheorem{corollary}[lemma]{Corollary}
\theoremstyle{remark}
\newtheorem{remark}[lemma]{Remark}
\newtheorem{convention}[lemma]{Convention}

\newcommand{\eps}{\varepsilon}
\newcommand{\Kfrak}{\mathfrak K}
\newcommand{\Kstar}{K^{*}}
\newcommand{\kxi}{\kappa_\xi}
\newcommand{\kbar}{\bar\kappa}
\newcommand{\lbar}{\bar\ell}
\newcommand{\epsbar}{\bar\eps}
\newcommand{\RII}{[R2]}

\title{Zone B: an analytic proof of the two-variable battle\\
{\large (replacement for the box certificate of Lemma \texttt{lem:zoneB-battle}
in \texttt{paper\_region2\_v2.tex})}}
\author{}
\date{2026-07-16}

\begin{document}
\maketitle

\noindent
\textbf{Scope.}  This note replaces the computational proof (the
$23$-box exhaustive-subdivision certificate \texttt{zoneB\_certifier.py})
of the Zone-B battle, Lemma~\texttt{lem:zoneB-battle} of
\texttt{paper\_region2\_v2.tex} (henceforth \RII), by a fully analytic
proof.  The \emph{statement} of the battle lemma is unchanged, so the proof
of Lemma~\texttt{lem:zoneB} in \RII{} (which combines the battle with
Lemma~\texttt{lem:Kmax}, the bounds
$\Pi\ge\underline\Pi$, $\Lambda\le\overline\Lambda$, and the certificate
\texttt{eq:cert-II}) goes through verbatim.  The only results quoted from
\RII{} are: the chart formulas of Lemma~\texttt{lem:chart}, and items (ii)
and (iv) of Lemma~\texttt{lem:Kmax}; everything else below is
self-contained.  Every displayed inequality is validated by
\texttt{validate\_zoneB.py} (dense float grids over the exact stated
domains, plus exact rational arithmetic for every table row, every
surrogate constant, and a family of corner spot checks).

\section{Setting and statement}

Throughout, $e\in[\tfrac1{60},\tfrac{2033}{10000}]$ and
\[
        \kxi(e)=\frac e{(1-e)^2},\qquad
        \kbar(e)=\min\Bigl(\frac{1-e}{1+e},\ \frac{1-3e}{6e}\Bigr),\qquad
        \kappa\in[\kxi(e),\kbar(e)]
\]
(the \emph{strip}; it is empty unless $\kxi\le\kbar$, and all statements are
vacuous at such $e$).  From the chart (Lemma~\texttt{lem:chart} of \RII):
\[
        x=x(e,\kappa)=\frac{1-e}{1+e+2\kappa e},\qquad
        u=1-x=\frac{2e(1+\kappa)}{1+e+2\kappa e},\qquad
        A=A(e,\kappa)=\frac{x(1+\kappa)}\kappa ,
\]
\[
        \Phi=\Phi(e,\kappa)=\frac{2(1-e)(1+\kappa)^2e}{\kappa(1+e+2\kappa e)^2},
        \qquad
        w(\nu)=\frac{\sqrt{\nu^2+4\nu}-\nu}2 ,
        \qquad
        \lambda=\ln(1/x).
\]
Since $\nu+2w(\nu)=\sqrt{\nu^2+4\nu}$, the quantity
$e^{-\Phi-2w(\Phi)}$ of \RII{} equals
\begin{equation}\label{eq:Kstar-def}
        \Kstar(e,\kappa):=\exp\Bigl(-\sqrt{\Phi^2+4\Phi}\Bigr).
\end{equation}
The battle lemma to be proved is
(Lemma~\texttt{lem:zoneB-battle} of \RII):
for all $e\in[\tfrac1{60},\tfrac{2033}{10000}]$ and
$\kappa\in[\kxi(e),\kbar(e)]$,
\begin{equation}\label{eq:battle}
        \underline\Pi(e,\kappa)\ \ge\
        \frac{\Kfrak(e,\kappa)}{x^2}+\overline\Lambda(e,\kappa),
\end{equation}
where, with $\lbar=\sqrt{2e/(1-e)}$,
$\underline\rho=\tfrac14(1-e^{-17.37(1+\kappa)e})$ and
$\epsbar=\min\bigl(\tfrac14,\;\frac e{4(1-e)^2\kappa(1+\underline\rho)}\bigr)$,
\[
        \underline\Pi
        =\sqrt{1-e}\;\frac{(1-\lbar)(1-\lbar^{14})(1-\epsbar)}{1+\lbar},
        \qquad
        \overline\Lambda
        =x^{13}\Bigl(\frac{2e}{1-e}\Bigr)^{13/2}\frac{(1-e)^2}{15e},
\]
and $\Kfrak=x^{14}(14-A)_+/15$ if the \emph{gate}
``$A<14$ and $15(14-A)\lambda\ge A+1$'' holds, while
$\Kfrak=\max\bigl(x^{14}(14-A)_+/15,\;\Kstar\bigr)$ otherwise.

Two explicit curves organise the proof:
\[
        \gamma(e):=\frac7{25}-e
        \qquad\text{(the \emph{gate curve})},
        \qquad
        G(e):=\sqrt{1-e}\;\frac{(1-\lbar)(1-\lbar^{14})}{1+\lbar},
\]
and the branch point of $\max(\gamma,\kxi)$ lies in
$(\tfrac{111}{1000},\tfrac18)$; we only use the rational comparison
$\kxi(\tfrac18)=\tfrac8{49}>\tfrac{31}{200}=\gamma(\tfrac18)$
(equivalently $8\cdot200>49\cdot31$).
The $e$-range is split at $e=\tfrac18$: for $e\le\tfrac18$ we work with the
curve $\kappa=\gamma(e)$, for $e\ge\tfrac18$ with the curve
$\kappa=\kxi(e)$.

\begin{convention}[safe decimals]\label{conv:safe}
Every terminating decimal below denotes the displayed exact rational
number.  In the verification tables, quantities that must be bounded
\emph{above} (marked~$\uparrow$) are exact rationals rounded up, and
quantities that must be bounded \emph{below} (marked~$\downarrow$) are
rounded down, so each table row is, as displayed, a finite exact rational
computation that certifies its inequality.  The script
\texttt{validate\_zoneB.py} re-verifies every row from the unrounded
rationals.
\end{convention}

\section{Monotonicity toolkit}

\begin{lemma}[elementary monotonicities]\label{lem:mono}
For $e\in(0,\tfrac13)$ and $\kappa>0$, with $D:=1+e+2\kappa e$:
\begin{enumerate}[label=\textup{(\roman*)},itemsep=1pt]
\item $x=(1-e)/D\in(0,1)$ is strictly decreasing in $e$ and in $\kappa$.
\item $u=1-x=2e(1+\kappa)/D$ is strictly increasing in $e$ and in $\kappa$:
$\partial_e u=2(1+\kappa)/D^2>0$ and $\partial_\kappa u=2e(1-e)/D^2>0$.
\item $A=(1-e)(1+\kappa)/(\kappa D)$ is strictly decreasing in $e$ and in
$\kappa$.
\item $\Phi$ is strictly decreasing in $\kappa$ on $\kappa\in(0,1]$;
and for $e\in[\tfrac1{60},\tfrac18]$, $\kappa\le\tfrac{79}{300}$, $\Phi$ is
strictly increasing in $e$.
\item $t\mapsto\exp(-\sqrt{t^2+4t})$ is strictly decreasing on $t\ge0$;
hence by \eqref{eq:Kstar-def}, $\Kstar$ is increasing in $\kappa$ on
$\kappa\in(0,1]$ and, on the domain of \textup{(iv)}, decreasing in $e$.
\item $G$ is positive and strictly decreasing on $(0,\tfrac13)$.
\item $\kxi$ is strictly increasing; $\gamma$ is strictly decreasing;
$\kxi/(4\gamma)$ is strictly increasing on $[\tfrac1{60},\tfrac18]$, with
maximal value $\kxi(\tfrac18)/(4\gamma(\tfrac18))=\tfrac{400}{1519}<1$.
\item $e\mapsto x(e,\gamma(e))=\dfrac{1-e}{1+\tfrac{39}{25}e-2e^2}$ is
strictly decreasing on $[\tfrac1{60},\tfrac18]$.
\item $e\mapsto x(e,\kxi(e))$ is strictly decreasing on
$[\tfrac18,\tfrac{2033}{10000}]$.
\item $\lambda=\ln(1/x)\ge u=1-x$ for $x\in(0,1]$.
\end{enumerate}
\end{lemma}

\begin{proof}
(i) In $\kappa$: $\partial_\kappa x=-2e(1-e)/D^2<0$.  In $e$: the numerator
$1-e$ decreases and $D=1+e(1+2\kappa)$ increases.
(ii) Direct differentiation as displayed.
(iii) In $e$: by (i), $x$ decreases and the factor $(1+\kappa)/\kappa$ is
constant.  In $\kappa$: with $N=1+\kappa$ and $\tilde D=\kappa D
=\kappa(1+e)+2\kappa^2e$,
\[
        N'\tilde D-N\tilde D'
        =\kappa(1+e)+2\kappa^2e-(1+\kappa)\bigl((1+e)+4\kappa e\bigr)
        =-(1+e)-4\kappa e-2\kappa^2e<0 .
\]
(iv) $\partial_\kappa\ln\Phi=\frac2{1+\kappa}-\frac1\kappa-\frac{4e}D
=\frac{\kappa-1}{\kappa(1+\kappa)}-\frac{4e}D\le-\frac{4e}D<0$ for
$\kappa\le1$.  Next,
$\partial_e\ln\Phi=\frac1e-\frac1{1-e}-\frac{2(1+2\kappa)}D
\ge\frac1e-\frac1{1-e}-2(1+2\kappa)$ (as $D\ge1$); for $e\le\tfrac18$ and
$\kappa\le\tfrac{79}{300}$ this is at least
$8-\tfrac87-2\bigl(1+\tfrac{158}{300}\bigr)=\tfrac{1997}{525}>0$.
(v) $t^2+4t$ is increasing on $t\ge0$; combine with (iv).
(vi) $\lbar=\sqrt{2e/(1-e)}$ is increasing in $e$ and $\lbar<1$ iff
$e<\tfrac13$; the three factors $\sqrt{1-e}$, $(1-\lbar)/(1+\lbar)$ and
$1-\lbar^{14}$ are positive and decreasing.
(vii) $\kxi$: numerator increases, denominator decreases.  The quotient of
a positive increasing by a positive decreasing function increases; its
value at $e=\tfrac18$ is $\frac{8/49}{4\cdot31/200}=\frac{400}{1519}$.
(viii) With $\kappa=\gamma(e)$ one gets $2\kappa e=\tfrac{14}{25}e-2e^2$,
so $D=1+\tfrac{39}{25}e-2e^2$, whose derivative $\tfrac{39}{25}-4e$ is
positive for $e\le\tfrac18<\tfrac{39}{100}$; the numerator $1-e$ decreases.
(ix) Composition of (i) with $\kxi$ increasing (vii).
(x) $-\ln x\ge1-x$.
\end{proof}

\begin{lemma}[the $\Lambda$-bound]\label{lem:lambda}
For all $(e,\kappa)$ in the strip,
\[
        \overline\Lambda(e,\kappa)\ \le\
        \Lambda_1(e):=\Bigl(\frac{1-e}{1+e}\Bigr)^{13}
        \Bigl(\frac{2e}{1-e}\Bigr)^{13/2}\frac{(1-e)^2}{15e}
        \ \le\ \widehat\Lambda:=\frac{13}{10^6}.
\]
\end{lemma}

\begin{proof}
The first inequality is $x\le(1-e)/(1+e)$ (Lemma~\ref{lem:mono}(i) with
$\kappa\downarrow0$), all other factors of $\overline\Lambda$ being
$\kappa$-free.  Collecting exponents,
\[
        \Lambda_1(e)=\frac{2^{13/2}}{15}\,
        \frac{e^{11/2}(1-e)^{17/2}}{(1+e)^{13}},
        \qquad
        (\ln\Lambda_1)'(e)=\frac{11}{2e}-\frac{17}{2(1-e)}-\frac{13}{1+e}.
\]
For $e\le e_1:=\tfrac{2033}{10000}$ each term is bounded by its value at
$e_1$ resp.\ by $13$:
\[
        (\ln\Lambda_1)'(e)\ \ge\ \frac{11}{2e_1}-\frac{17}{2(1-e_1)}-13
        =\frac{55000}{2033}-\frac{85000}{7967}-13>3.38>0
\]
(exact rational comparison: $55000\cdot7967-85000\cdot2033-13\cdot2033\cdot
7967=54{,}820{,}157>0$).  Hence $\Lambda_1$ is increasing on $(0,e_1]$ and
\[
        \Lambda_1(e)\le\Lambda_1(e_1)
        =\frac{2^6}{15}\,\sqrt{2e_1(1-e_1)}\;
        \frac{e_1^5(1-e_1)^8}{(1+e_1)^{13}}
        \ \le\ \frac{64}{15}\cdot\frac{14229}{25000}\cdot
        \frac{e_1^5(1-e_1)^8}{(1+e_1)^{13}}
        \ \le\ \frac{124}{10^7}\ <\ \widehat\Lambda,
\]
using $\bigl(\tfrac{14229}{25000}\bigr)^2\ge2e_1(1-e_1)
=\tfrac{32393822}{10^8}$; the final two comparisons are exact rational
evaluations (script, \S1--2).
\end{proof}

\begin{lemma}[envelope and the $\Kfrak$-bounds]\label{lem:envelope}
At every strip point:
\begin{enumerate}[label=\textup{(\alph*)},itemsep=1pt]
\item $\underline\Pi\ \ge\ G(e)\,(1-\epsbar)$ with
$\epsbar\le\min\bigl(\tfrac14,\ \tfrac{\kxi(e)}{4\kappa}\bigr)$;
in particular $\underline\Pi\ge\tfrac34G(e)$.
\item $\Kfrak\le\Kstar(e,\kappa)$ (in both gate branches); consequently the
$\max(\cdot,\cdot)$ in the definition of $\Kfrak$ is redundant.
\item If the gate holds at $(e,\kappa)$ then
$\Kfrak\le\tfrac{14}{15}x^{14}$.
\end{enumerate}
\end{lemma}

\begin{proof}
(a) $\underline\Pi=G(e)(1-\epsbar)$ by definition, and
$\epsbar\le\frac e{4(1-e)^2\kappa(1+\underline\rho)}
\le\frac e{4(1-e)^2\kappa}=\frac{\kxi(e)}{4\kappa}$ because
$\underline\rho\ge0$; also $\epsbar\le\tfrac14$ by its definition as a
minimum, and $\kappa\ge\kxi$ gives again $\epsbar\le\tfrac14$.

(b) Write $K(n)=x^n(n-A)/(n+1)$.  If $A\ge14$ then
$x^{14}(14-A)_+/15=0\le\Kstar$.  If $A<14$, then by
Lemma~\texttt{lem:Kmax}(ii) of \RII{} the function $K$ is log-concave on
$(A,\infty)$ with an interior stationary point $n_*$ (the unique root of
$\lambda=(A+1)/((n-A)(n+1))$, which exists because the right side decreases
from $+\infty$ to $0$), so $K(14)\le K(n_*)$; and by
Lemma~\texttt{lem:Kmax}(iv), $K(n_*)\le e^{-\Phi-2w(\Phi)}=\Kstar$.  Thus
$x^{14}(14-A)_+/15\le\Kstar$ always.  On the gate branch
$\Kfrak=x^{14}(14-A)_+/15\le\Kstar$; on the other branch
$\Kfrak=\max\bigl(x^{14}(14-A)_+/15,\Kstar\bigr)=\Kstar$.

(c) On the gate branch, $(14-A)_+\le14$.
\end{proof}

\section{The gate on the two curves}

\begin{lemma}[$u$-gate]\label{lem:ugate}
Say the \emph{$u$-gate} holds at $(e,\kappa)$ if $A<14$ and
$15(14-A)u\ge A+1$.  Then:
\begin{enumerate}[label=\textup{(\alph*)},itemsep=1pt]
\item ($u$-gate $\Rightarrow$ gate)  If the $u$-gate holds at $(e,\kappa)$,
so does the gate of \eqref{eq:battle}, since
$\lambda\ge u$ (Lemma~\ref{lem:mono}(x)) and $14-A>0$ give
$15(14-A)\lambda\ge15(14-A)u\ge A+1$.
\item (upward extension)  If the $u$-gate holds at $(e,\kappa_0)$, it holds
at $(e,\kappa)$ for every $\kappa\ge\kappa_0$: by
Lemma~\ref{lem:mono}(ii),(iii), $A(e,\kappa)\le A(e,\kappa_0)<14$ and
$u(e,\kappa)\ge u(e,\kappa_0)$, so
$15\bigl(14-A(e,\kappa)\bigr)u(e,\kappa)
\ge15\bigl(14-A(e,\kappa_0)\bigr)u(e,\kappa_0)\ge A(e,\kappa_0)+1\ge
A(e,\kappa)+1$. \qedhere
\end{enumerate}
\end{lemma}

\begin{proposition}[gate on the curves]\label{prop:gate}
\begin{enumerate}[label=\textup{(\alph*)},itemsep=1pt]
\item For every $e\in[\tfrac1{60},\tfrac18]$ the $u$-gate holds at
$(e,\gamma(e))$, hence the gate holds at $(e,\kappa)$ for all
$\kappa\ge\gamma(e)$.
\item For every $e\in[\tfrac18,\tfrac{2033}{10000}]$ the $u$-gate holds at
$(e,\kxi(e))$, hence the gate holds at every strip point with
$e\ge\tfrac18$.
\end{enumerate}
\end{proposition}

\begin{proof}
(a) Fix a row $[a,b]$ of Table~\ref{tab:gate} (the three rows tile
$[\tfrac1{60},\tfrac18]$) and let $e\in[a,b]$, $\kappa=\gamma(e)\in
[\gamma(b),\gamma(a)]$.  By Lemma~\ref{lem:mono}(i),(iii),
\[
        A(e,\gamma(e))=x(e,\gamma(e))\Bigl(1+\frac1{\gamma(e)}\Bigr)
        \ \le\ x\bigl(a,\gamma(b)\bigr)\Bigl(1+\frac1{\gamma(b)}\Bigr)
        =:\overline A,
\]
since $x$ decreases in both arguments and $1+1/\gamma(e)\le1+1/\gamma(b)$;
and by Lemma~\ref{lem:mono}(ii),
$u(e,\gamma(e))\ge u(a,\gamma(b))=:\underline u$.
The row certifies $\overline A<14$ and
$15(14-\overline A)\,\underline u-(\overline A+1)>0$; monotonicity then
gives $15(14-A)u\ge15(14-\overline A)\underline u\ge\overline A+1\ge A+1$,
i.e.\ the $u$-gate at $(e,\gamma(e))$.  Lemma~\ref{lem:ugate} extends it
upward and converts it to the gate.

(b) Same argument on the single row $[\tfrac18,\tfrac{2033}{10000}]$: by
Lemma~\ref{lem:mono}(ii),(iii),(vii), for $e\ge a=\tfrac18$,
\[
        A(e,\kxi(e))\le A(e,\kxi(a))\le A(a,\kxi(a))=:\overline A,
        \qquad
        u(e,\kxi(e))\ge u(a,\kxi(a))=:\underline u ,
\]
and the row of Table~\ref{tab:gate} certifies the $u$-gate.  Every strip
point with $e\ge\tfrac18$ has $\kappa\ge\kxi(e)$, so
Lemma~\ref{lem:ugate}(b),(a) applies.
\end{proof}

\begin{table}[h]\centering
\caption{Gate rows.  $\overline A{=}x(a,\kappa_-)(1{+}1/\kappa_-)\uparrow$,
$\underline u{=}u(a,\kappa_-)\downarrow$ with $\kappa_-=\gamma(b)$ (part a)
resp.\ $\kappa_-=\kxi(a)$ (part b); slack
$=15(14-\overline A)\underline u-(\overline A+1)\downarrow$.  All entries
exact rationals (Convention~\ref{conv:safe}).}\label{tab:gate}
\begin{tabular}{lcccc}
\toprule
$[a,b]$ & $\kappa_-$ & $\overline A\uparrow$ & $\underline u\downarrow$ & slack$\downarrow$\\
\midrule
\multicolumn{5}{l}{(a) curve $\kappa=\gamma(e)$ on $[\tfrac1{60},\tfrac18]$}\\
$[1/60,\,21/1000]$ & $259/1000$ & $4.6621$ & $0.0409$ & $0.06$ \\
$[21/1000,\,8/125]$ & $27/125$ & $5.3506$ & $0.0495$ & $0.07$ \\
$[8/125,\,1/8]$ & $31/200$ & $6.4352$ & $0.1364$ & $8.04$ \\
\midrule
\multicolumn{5}{l}{(b) curve $\kappa=\kxi(e)$ on $[\tfrac18,\tfrac{2033}{10000}]$}\\
$[1/8,\,2033/10000]$ & $8/49$ & $5.3477$ & $0.2494$ & $26.02$ \\
\bottomrule
\end{tabular}
\end{table}

\section{The two one-dimensional inequalities}

For a rational $b$ let $c_\ell(b)$, $c_s(b)$ denote the $5$-digit decimals
of Table~\ref{tab:surr} with
\begin{equation}\label{eq:surr}
        c_\ell^2\ \ge\ \frac{2b}{1-b},
        \qquad
        c_s^2\ \le\ 1-b
        \qquad(\text{exact rational comparisons}),
\end{equation}
so that, by Lemma~\ref{lem:mono}(vi) and since
$\lbar(b)^{14}=\bigl(\tfrac{2b}{1-b}\bigr)^7$ is rational,
\begin{equation}\label{eq:Glo}
        G(e)\ \ge\ G(b)\ \ge\
        G_\downarrow(b):=c_s\,
        \frac{\bigl(1-c_\ell\bigr)\Bigl(1-\bigl(\tfrac{2b}{1-b}\bigr)^{7}\Bigr)}
        {1+c_\ell}
        \qquad\text{for } e\le b .
\end{equation}

\begin{proposition}[Case-S inequality]\label{prop:caseS}
For every $e\in[\tfrac1{60},\tfrac18]$,
\begin{equation}\label{eq:caseS}
        \tfrac34\,G(e)\ \ge\
        \frac{\Kstar\bigl(e,\gamma(e)\bigr)}{x\bigl(e,\gamma(e)\bigr)^2}
        +\widehat\Lambda .
\end{equation}
\end{proposition}

\begin{proof}
Fix the row $[a,b]$ of Table~\ref{tab:caseS} containing $e$ (the seven rows
tile $[\tfrac1{60},\tfrac18]$).  Since $\gamma(e)\le\gamma(a)\le
\gamma(\tfrac1{60})=\tfrac{79}{300}$ and by
Lemma~\ref{lem:mono}(iv) (both parts, on the stated domains),
\[
        \Phi\bigl(e,\gamma(e)\bigr)\ \ge\ \Phi\bigl(e,\gamma(a)\bigr)
        \ \ge\ \Phi\bigl(a,\gamma(a)\bigr)=:\Phi_a ,
\]
so by Lemma~\ref{lem:mono}(v) and $e^t\ge1+t+\tfrac{t^2}2+\tfrac{t^3}6$
($t\ge0$),
\[
        \Kstar\bigl(e,\gamma(e)\bigr)
        \ \le\ \exp\Bigl(-\sqrt{\Phi_a^2+4\Phi_a}\Bigr)
        \ \le\ e^{-r}
        \ \le\ \frac1{1+r+\frac{r^2}2+\frac{r^3}6}=:K^{\uparrow},
\]
where $r=r(a)$ is the rational surrogate of Table~\ref{tab:caseS} with
$r^2\le\Phi_a^2+4\Phi_a$ (exact comparison).  By
Lemma~\ref{lem:mono}(viii), $x(e,\gamma(e))^2\ge x(b,\gamma(b))^2$, and by
\eqref{eq:Glo}, $G(e)\ge G_\downarrow(b)$.  The row certifies
\[
        \tfrac34\,G_\downarrow(b)\ \ge\
        \frac{K^{\uparrow}}{x(b,\gamma(b))^2}+\widehat\Lambda,
\]
which combined with the three monotone bounds yields \eqref{eq:caseS}.
\end{proof}

\begin{table}[h]\centering
\caption{Case-S rows.  $r=r(a)$ satisfies $r^2\le\Phi_a(\Phi_a+4)$,
$\Phi_a=\Phi(a,\gamma(a))$;
$K^\uparrow=1/(1{+}r{+}\tfrac{r^2}2{+}\tfrac{r^3}6)\uparrow$;
$X_\downarrow=x(b,\gamma(b))^2\downarrow$;
margin $=\tfrac34G_\downarrow(b)-K^\uparrow/X_\downarrow-\widehat\Lambda
\downarrow$.}\label{tab:caseS}
\begin{tabular}{lccccc}
\toprule
$[a,b]$ & $r$ & $K^\uparrow$ & $X_\downarrow$ & $G_\downarrow(b)\downarrow$ & margin$\downarrow$\\
\midrule
$[1/60,\,23/1000]$ & $0.88959$ & $0.41622$ & $0.89136$ & $0.63595$ & $0.0100$ \\
$[23/1000,\,33/1000]$ & $1.04817$ & $0.35850$ & $0.84928$ & $0.57596$ & $0.0098$ \\
$[33/1000,\,6/125]$ & $1.26209$ & $0.29468$ & $0.79119$ & $0.50537$ & $0.0065$ \\
$[6/125,\,67/1000]$ & $1.53635$ & $0.23144$ & $0.72527$ & $0.43499$ & $0.0071$ \\
$[67/1000,\,89/1000]$ & $1.84283$ & $0.17909$ & $0.65808$ & $0.36930$ & $0.0048$ \\
$[89/1000,\,111/1000]$ & $2.17522$ & $0.13781$ & $0.59913$ & $0.31450$ & $0.0058$ \\
$[111/1000,\,1/8]$ & $2.51056$ & $0.10754$ & $0.56532$ & $0.28369$ & $0.0225$ \\
\bottomrule
\end{tabular}
\end{table}

\begin{proposition}[Case-L inequalities]\label{prop:caseL}
\begin{enumerate}[label=\textup{(\alph*)},itemsep=1pt]
\item For every $e\in[\tfrac1{60},\tfrac18]$,
\begin{equation}\label{eq:caseLg}
        G(e)\Bigl(1-\frac{\kxi(e)}{4\gamma(e)}\Bigr)
        \ \ge\ \tfrac{14}{15}\,x\bigl(e,\gamma(e)\bigr)^{12}+\widehat\Lambda .
\end{equation}
\item For every $e\in[\tfrac18,\tfrac{2033}{10000}]$,
\begin{equation}\label{eq:caseLx}
        \tfrac34\,G(e)
        \ \ge\ \tfrac{14}{15}\,x\bigl(e,\kxi(e)\bigr)^{12}+\widehat\Lambda .
\end{equation}
\end{enumerate}
\end{proposition}

\begin{proof}
(a) Fix the row $[a,b]$ of Table~\ref{tab:caseL}(a) containing $e$.  By
Lemma~\ref{lem:mono}(vi),(vii) the left side of \eqref{eq:caseLg} is a
product of two positive decreasing functions of $e$ (positivity of the
second factor: $\kxi/(4\gamma)\le\tfrac{400}{1519}<1$ on
$[\tfrac1{60},\tfrac18]$ by (vii)), so it is at least
$G_\downarrow(b)\bigl(1-\kxi(b)/(4\gamma(b))\bigr)=:L_\downarrow$; by
Lemma~\ref{lem:mono}(viii) the right side is at most
$\tfrac{14}{15}x(a,\gamma(a))^{12}+\widehat\Lambda=:R_\uparrow$ (the power
$x(a,\gamma(a))^{12}$ is an exact rational).  The row certifies
$L_\downarrow\ge R_\uparrow$.

(b) Same with the single row of Table~\ref{tab:caseL}(b), using
Lemma~\ref{lem:mono}(vi) for the left side and (ix) for the right.
\end{proof}

\begin{table}[h]\centering
\caption{Case-L rows.  (a): $L_\downarrow=G_\downarrow(b)\bigl(1-
\kxi(b)/(4\gamma(b))\bigr)\downarrow$,
$R_\uparrow=\tfrac{14}{15}x(a,\gamma(a))^{12}\uparrow$.
(b): $L_\downarrow=\tfrac34G_\downarrow(b)\downarrow$,
$R_\uparrow=\tfrac{14}{15}x(a,\kxi(a))^{12}\uparrow$.
Margin $=L_\downarrow-R_\uparrow-\widehat\Lambda\downarrow$.}\label{tab:caseL}
\begin{tabular}{lccc}
\toprule
$[a,b]$ & $L_\downarrow$ & $R_\uparrow$ & margin$\downarrow$\\
\midrule
\multicolumn{4}{l}{(a) curve $\kappa=\gamma(e)$ on $[\tfrac1{60},\tfrac18]$}\\
$[1/60,\,29/1000]$ & $0.57998$ & $0.56429$ & $0.0156$ \\
$[29/1000,\,13/200]$ & $0.40350$ & $0.39307$ & $0.0104$ \\
$[13/200,\,1/8]$ & $0.20898$ & $0.14341$ & $0.0655$ \\
\midrule
\multicolumn{4}{l}{(b) curve $\kappa=\kxi(e)$ on $[\tfrac18,\tfrac{2033}{10000}]$}\\
$[1/8,\,2033/10000]$ & $0.11051$ & $0.02983$ & $0.0806$ \\
\bottomrule
\end{tabular}
\end{table}

\begin{table}[h]\centering
\caption{Surrogates for \eqref{eq:surr} (exact $5$-digit decimals;
each satisfies the displayed exact rational comparison).}\label{tab:surr}
\begin{tabular}{lcc@{\qquad}lcc}
\toprule
$b$ & $c_\ell$ & $c_s$ & $b$ & $c_\ell$ & $c_s$\\
\midrule
$23/1000$ & $0.21699$ & $0.98843$ & $89/1000$ & $0.44203$ & $0.95446$\\
$29/1000$ & $0.24441$ & $0.98539$ & $111/1000$ & $0.49972$ & $0.94286$\\
$33/1000$ & $0.26126$ & $0.98336$ & $1/8$ & $0.53453$ & $0.93541$\\
$6/125$ & $0.31756$ & $0.97570$ & $2033/10000$ & $0.71440$ & $0.89258$\\
$13/200$ & $0.37288$ & $0.96695$ & & &\\
$67/1000$ & $0.37898$ & $0.96591$ & & &\\
\bottomrule
\end{tabular}
\end{table}

\section{Proof of the battle}

\begin{proof}[Proof of \eqref{eq:battle} \textup{(Lemma
\texttt{lem:zoneB-battle} of \RII)}]
Let $(e,\kappa)$ be a strip point.  By Lemma~\ref{lem:lambda},
$\overline\Lambda\le\widehat\Lambda$ throughout, so it suffices to prove
$\underline\Pi\ge\Kfrak/x^2+\widehat\Lambda$.  We distinguish three cases.

\medskip\noindent
\emph{Case S: $e\le\tfrac18$ and $\kappa\le\gamma(e)$.}
By Lemma~\ref{lem:envelope}(a),(b),
\[
        \underline\Pi\ \ge\ \tfrac34G(e),
        \qquad
        \frac{\Kfrak}{x^2}\ \le\ \frac{\Kstar(e,\kappa)}{x(e,\kappa)^2}.
\]
Since $\kappa\le\gamma(e)\le\tfrac{79}{300}<1$, Lemma~\ref{lem:mono}(iv),(v)
give $\Kstar(e,\kappa)\le\Kstar(e,\gamma(e))$, and Lemma~\ref{lem:mono}(i)
gives $x(e,\kappa)\ge x(e,\gamma(e))$; hence
$\Kfrak/x^2\le\Kstar(e,\gamma(e))/x(e,\gamma(e))^2$, and
Proposition~\ref{prop:caseS} closes this case.

\medskip\noindent
\emph{Case L-$\gamma$: $e\le\tfrac18$ and $\kappa\ge\gamma(e)$.}
By Proposition~\ref{prop:gate}(a) the gate holds at $(e,\kappa)$, so
Lemma~\ref{lem:envelope}(c) and Lemma~\ref{lem:mono}(i) give
\[
        \frac{\Kfrak}{x^2}\ \le\ \tfrac{14}{15}\,x(e,\kappa)^{12}
        \ \le\ \tfrac{14}{15}\,x\bigl(e,\gamma(e)\bigr)^{12}.
\]
By Lemma~\ref{lem:envelope}(a) and $\kappa\ge\gamma(e)$,
\[
        \underline\Pi\ \ge\ G(e)\Bigl(1-\frac{\kxi(e)}{4\kappa}\Bigr)
        \ \ge\ G(e)\Bigl(1-\frac{\kxi(e)}{4\gamma(e)}\Bigr),
\]
and Proposition~\ref{prop:caseL}(a) closes this case.

\medskip\noindent
\emph{Case L-$\xi$: $e\ge\tfrac18$.}
Here $\kappa\ge\kxi(e)$, so by Proposition~\ref{prop:gate}(b) the gate
holds; as above
$\Kfrak/x^2\le\tfrac{14}{15}x(e,\kappa)^{12}
\le\tfrac{14}{15}x(e,\kxi(e))^{12}$, while
$\underline\Pi\ge\tfrac34G(e)$ by Lemma~\ref{lem:envelope}(a).
Proposition~\ref{prop:caseL}(b) closes this case.

\medskip
The three cases cover the strip: for $e\le\tfrac18$ every $\kappa$
satisfies $\kappa\le\gamma(e)$ or $\kappa\ge\gamma(e)$, and $e\ge\tfrac18$
is Case L-$\xi$.
\end{proof}

\section{Remarks}

\begin{remark}[the $\underline\rho$-term is not needed]\label{rmk:rho}
The proof used only $\underline\rho\ge0$
(Lemma~\ref{lem:envelope}(a)).  Hence the battle \eqref{eq:battle} holds
with $\epsbar$ replaced by the larger
$\min\bigl(\tfrac14,\kxi(e)/(4\kappa)\bigr)$, i.e.\ the
$17.37$-exponential can be deleted from the statement of
Lemma~\texttt{lem:zoneB-battle}, and with it the rational comparison
$28(1+\kappa)e/(1+3e)\ge17.37(1+\kappa)e$ in the proof of
Lemma~\texttt{lem:zoneB} of \RII{} becomes unnecessary (only
$\rho\ge0$ is used there).
\end{remark}

\begin{remark}[the $\max$ is redundant]
By Lemma~\ref{lem:envelope}(b), on the gate-failure branch
$\Kfrak=\Kstar$; the paper may state the dichotomy as: gate
$\Rightarrow\Kfrak=x^{14}(14-A)_+/15$, otherwise $\Kfrak=\Kstar$.
\end{remark}

\begin{remark}[margins]
The battle \eqref{eq:battle} in fact holds with pointwise margin
$\ge0.116$ on the whole strip (minimum at the pinch corner
$e\approx0.2029$, $\kappa=\kxi(e)$, where $\kxi$ meets $\kappa_q$); the
analytic chain above retains a margin $\ge0.0048$ in the tightest table
row (Case S, $[67/1000,89/1000]$) and $\ge0.06$ in the tightest gate row.
The proof uses $15$ table rows in total.
\end{remark}

\begin{remark}[verification protocol]
\texttt{validate\_zoneB.py} checks, and exits $0$ only if all pass:
(1) exact rational verification of every row of
Tables~\ref{tab:gate}--\ref{tab:surr}, including every surrogate comparison
\eqref{eq:surr} and $r^2\le\Phi_a(\Phi_a+4)$;
(2) the exact rational comparisons quoted in Lemma~\ref{lem:lambda} and
Lemma~\ref{lem:mono}(iv),(vii);
(3) every monotonicity claim of Lemma~\ref{lem:mono} on a dense grid over
its exact stated domain;
(4) an end-to-end pointwise replay of the three-case assembly on a
$1201\times1201$ grid over the strip (each intermediate inequality
checked pointwise, including $\Kfrak\le\Kstar$ and the direct battle
margin);
(5) exact-\texttt{Fraction} spot checks of \eqref{eq:battle} at $25$
corner/extreme rational points via safe-direction brackets, independent of
the tables.
\end{remark}

\end{document}
